\anexo{Transcripción de la entrevista a Luciana Jacubovich}

\textbf{Entrevistada:} Jacubovich Luciana (responsable de seguridad y gestión de accesos, grupo económico financiero)

\textbf{Entrevistadores:} Muntaabski Federico, Mociulsky Santiago Bernardo

\textbf{Nota de transcripción:} el archivo original identificaba a los entrevistadores como "Speaker 2", "Speaker 3" y "Speaker 4" (probable artefacto de diarización automática que separó a los mismos dos entrevistadores en más de una etiqueta). Se asignaron los nombres reales según el patrón usado en las demás entrevistas (Federico pregunta primero, Santiago repregunta); \textbf{conviene revisar manualmente} si alguna intervención puntual quedó mal atribuida entre Federico y Santiago.

\begin{description}[leftmargin=0cm, labelsep=0.5cm]
    \item[\textbf{JACUBOVICH LUCIANA:}] Yo les cuento, yo trabajo en lo que es seguridad hace más de 20 años. Empecé como analista de seguridad en una empresa de seguros y hoy en día lidero un equipo de 7 personas dentro de lo que es el área de operaciones de seguridad y accesos, gestión de accesos de un grupo económico que se llama Grupo ST. El grupo es un grupo financiero, o sea, que está conformado por una aseguradora, un banco, una SGR, o sea, una compañía que brinda leasing, o sea, es un múltiples compañías que conforman este grupo financiero. Y bueno, yo estoy a cargo de todo lo que dentro del área de lo que sería seguridad, de todo lo que es seguridad en operaciones y gestión de accesos. O sea, sería una gran interesada del producto de ustedes para obtenerlo. Así que, nada, está buenísima la idea, me parece buenísima. Ahí como para darle una utilidad y poder generar una acción más proactiva para la gestión de los accesos de los usuarios y mitigar riesgos. Bárbaro, bárbaro.

    \item[\textbf{MUNTAABSKI FEDERICO:}] Perfecto, porque eso, o sea, nos viene perfecto también entender un poco cuál es el trasfondo que tenés vos de carrera, pues también las preguntas nos permite también encaminarlas un poco por eso. Buenísimo. Bueno, ¿te molesta si empezamos la entrevista? Tranquilo, disparen. Y como consultita.

    \item[\textbf{MOCIULSKY SANTIAGO BERNARDO:}] Más que nada por el tema de transcripción, la minuta después, ¿te molestas y grabamos? Para nada, chicos, graben. Genial, gracias. Bueno, ¿pones vos, Santi? Sí, ahí te grabamos. Ok, ahí... Ah, perfecto. Ahí está, ahora ya apareció. Ahí está. Bueno, ahora vamos a arrancar la entrevista. Dale.

    \item[\textbf{MUNTAABSKI FEDERICO:}] Vamos a ir con lo que es la primera pregunta. La primera pregunta sería, ¿cómo describirías el mayor desafío de seguridad que hoy ves relacionado a tu organización o empresa? El mayor desafío hoy es la inteligencia artificial. No solamente saber cómo usarla, sino también cómo lograr.

    \item[\textbf{JACUBOVICH LUCIANA:}] Que los usuarios que la usan no expongan información de las compañías y de nuestros clientes, en definitiva. El objeto final del área de seguridad nuestro siempre es cuidar la información de nuestros clientes. Entonces, lo fundamental es eso, que no se termine exponiendo ninguna información. como activo de la compañía, como información de clientes nuestros, a través de una buena intención del uso de la inteligencia artificial para mejorar algún proceso o algo. Entonces el mayor desafío hoy está en poder potenciar a la compañía permitiendo el uso de esa inteligencia artificial, pero en forma limitada para que esto no suceda. La gran disyuntiva está ahí porque hay una gran presión de los accionistas para poder hacer uso de estas inteligencias artificiales y potenciar a la compañía, pero a su vez también nosotros tenemos que hablar por la información de nuestros clientes y que no sea expuesta. Entiendo, perfecto. Claro, ahí está la acción entre medir los beneficios que traen estas herramientas de inteligencia artificial. contra los riesgos que presenta en la privacidad y todo lo relacionado a los datos. Sí, después también está la arista, no lo nombré, pero también es importante, los atacantes también utilizan esta inteligencia artificial y entonces los ataques... Se mueven mucho más rápido, se movilizan mucho más rápido que en otras épocas o simulan mucho más a un humano que lo que lo hacían antes. Antes uno era mucho más fácil de detectar cuando tenías algún bot o algo atrás que te estaba queriendo atacar a través de alguna herramienta que estaba utilizando algún atacante. Hoy ya eso es más difícil de detectar, dado la... al uso de la inteligencia artificial atrás y la propia inteligencia ya va aprendiendo y va cambiando su comportamiento, entonces eso también hace que sea más desafiante el poder cuidar los activos de la compañía de esa forma. O sea, si me preguntas hoy cuál es tu desafío de la inteligencia artificial en todos sus aristas, o sea, desde el uso hasta el cuidado y hasta el poder

    \item[\textbf{MUNTAABSKI FEDERICO:}] protegernos de los atacantes que hacen uso de esa inteligencia artificial. Claro. Perfecto. Bueno, pasamos a la siguiente pregunta, que sería, ¿qué tipos de incidentes de seguridad relacionados con usuarios son los que más han tenido impacto en tu empresa o organización en las que has trabajado en los últimos años? Sí.

    \item[\textbf{JACUBOVICH LUCIANA:}] Se me ocurren dos, o sea, como principales, filtración de datos inintencionadamente, digamos, sin intención de algún usuario que mandó algún archivo con algún adjunto o algo, alguna dirección errónea o alguien que no lo tenía que mandar. Ese es uno de los incidentes más comunes, digamos, con mayor impacto. Y por otro lado, los phishings. O sea, parece ya viejo, pero sigue estando presente todos los días y evoluciona con esto también de la inteligencia artificial. Cada vez es más artesanal, más dirigido a cada persona, pero la gente sigue cayendo en el phishing, entonces también trabajamos muy duro en concientizar a los usuarios en cuanto a eso.

    \item[\textbf{MOCIULSKY SANTIAGO BERNARDO:}] Claro, perfecto. O sea, el phishing es una de las principales problemáticas que se presentan. Sí. Perfecto. Bueno, si querés, hago yo un poco para variar. Después te quería consultar más orientado a la gestión de acceso, que vos dijiste que estabas laburando en eso. Con respecto a la supervisación o el monitoreo de usuarios, ¿ustedes qué nivel de acciones pueden llegar a ver?

    \item[\textbf{JACUBOVICH LUCIANA:}] ¿Qué tanto pueden ver de la actividad del día a día del usuario? A ver, tenés diferentes capas, ¿no? O sea, por lo general uno tenés como capas o aristas, digamos, de lo que hace un usuario. Podés verlo desde que se conecta a nivel red, ¿sí? ¿A qué webs? está accediendo o qué aplicaciones tiene instalada en su máquina, un usuario, a qué aplicaciones, después tenés las actividades que realiza dentro de un aplicativo, si el aplicativo tiene un logueo, depende del nivel al que loguea la aplicación, eso va a depender de cada aplicación. Nosotros tenemos ciertos estándares, pero muchas veces vienen aplicaciones ya dadas que no tienen. que solamente tienen un login básico, de un login o un logout, y ya está, te quedaste ahí, y no sabes qué acciones realizan dentro del aplicativo. Cuando son aplicaciones críticas de negocio, nosotros solicitamos que tengan ese login de, bueno, realizó esta acción crítica dentro de la aplicación, ¿no? No sé, por decir algo relacionado a seguros, modificó el monto de la prima de una póliza, ¿sí? O el monto asegurado de una póliza. Entonces, esa acción es crítica, bueno. Se tiene que quedar registrado que esa persona realizó esa acción. Entonces, va a depender, en ese caso, del aplicativo que tanto lo intenia. Pero, como te decías, como uno va a depender el tipo de usuario, también cuánta profundidad va a tener. Porque capaz un usuario que no accede a ninguna aplicación, sino que meramente utiliza, no sé, su correo electrónico y navega por internet.

    \item[\textbf{MOCIULSKY SANTIAGO BERNARDO:}] No vas a tener mucho nivel de detalle, más que eso, pero va a depender un poco de la actividad y los privilegios también que tenga ese usuario. Ok, y por ejemplo, los usuarios que tienen privilegios, por ejemplo, para hacer la base de datos o datos críticos, ¿será este un seguimiento de actividad que realizan o cómo es ese chequeo? Se loguea, queda logueado, queda registrado.

    \item[\textbf{JACUBOVICH LUCIANA:}] Pero no quiere decir que se haga un seguimiento, ¿sí? No se revisa ese log en forma diaria, sino que lo que se hace es como el control inverso. Si hay alguien que está realizando una tarea crítica que no está dentro del grupo de usuarios que está permitido, ahí se genera una alerta. ¿Sí? Es como la situación inversa. Porque tenés ciertas personas que ya tienen esa actividad permitida. Yo de esas personas sí voy a tener un registro de qué es lo que están haciendo, por una cuestión de control, si requiero alguna auditoría posterior a eso, pero de la situación como más reactiva, preventiva, lo que voy a tener es un control de una alerta si hay una actividad de, por ejemplo, una persona que tiene privilegios de administración de base de datos. pero que no está dentro de las personas que yo sé que son las administradoras de base de datos. ¿Sí? Que por algún motivo tengo algún usuario que ganó esos privilegios y puedo detectarlo fácilmente porque está ejecutando tareas que no son propias de su usuario. ¿No? Ok. Genial. Sería como medio un control por oposición de alguna forma. Ok.

    \item[\textbf{MOCIULSKY SANTIAGO BERNARDO:}] Y con respecto a esto que vos nombraste, los registros, lo de auditoría, ¿qué herramientas se utiliza para eso? Tenemos un SIEM, que es el de Azure, el Sentinel, donde se colectan ahí todos los eventos. Después cada plataforma, cada herramienta tiene su propio login, pero lo que hacemos es recolectarlos después todos. en este aplicativo que es el Sentinel entonces ahí lo que puedes empezar a hacer son las correlaciones estas de eventos genial eso se entendió perfecto ahora pasamos a otra pregunta que es principalmente ¿cuánto tiempo pasa de por ejemplo que ocurre una actividad que es sospechosa dentro de un usuario hasta que verdaderamente el equipo lo detecte y lo soluciona? ¿se entiende? más o menos

    \item[\textbf{JACUBOVICH LUCIANA:}] Sí, va a depender de la criticidad de ese evento, ¿no? Si hay una actividad sospechosa probablemente esté siendo alertada y eso es, te diría que casi inmediato de cuando ocurre esa actividad sospechosa. Ahora, después hay todo un proceso de investigación, que eso sí ya puede llevar más tiempo, pues, no sé, desde... Una hora, hasta cinco horas, depende de qué tanto te cueste rastrear eso. La idea es justamente tener todos los eventos dentro de la consola para poder hacer ese análisis, pero muchas veces requiere poder contactar al usuario, al equipo del usuario, secuestrar de alguna forma al equipo del usuario, tomar posesión de ese equipo para poder hacer una investigación más profunda por si el equipo estuvo... tomado por el incidente entonces va a depender pero desde el momento en que se detecta que es casi al unísono de que sucede porque nada tenemos las alertas y tenemos la consola generando esos eventos después el tiempo de investigación va a depender del tipo de incidente

    \item[\textbf{MUNTAABSKI FEDERICO:}] Una pregunta medio relacionada por ese estilo de los incidentes. En el caso de que ustedes detecten una actividad sospechosa de algún usuario, ¿cómo sería el proceso desde que se detecta hasta que se llevan las acciones? Perdón, me te consulto. Una pregunta que viene un poco antes de esa es, ¿usted tiene un equipo SOC? supongo, dentro de la organización. Sí. Chiquito, pero sí. Todavía somos medio APL en ese sentido, pero sí. Estamos evolucionando, vamos de a poco evolucionando. Así que no, la pregunta era, primero, si tenían equipo SOC y después cómo es ese proceso desde que se detecta la amenaza hasta el proceso en el cual

    \item[\textbf{JACUBOVICH LUCIANA:}] Se finaliza la amenaza ya sea bloqueando o haciendo la acción correspondiente. Primero, aclaración es, yo tengo un evento. No sé si es un incidente o no es un incidente, que es un poco lo que les decía antes. Yo tengo esa alerta, a mí me llega esa alerta y no sé en qué deriva. Por eso en cantidad de horas no les puedo decir, el promedio es tanto. Si querés te digo, el promedio son dos horas de análisis, ¿sí? Por lo general, hasta que uno dice, fue un incidente, no fue un incidente, está contenido, no está contenido, o sea, en el promedio son dos horas, ¿sí? Pero hay veces que son diez minutos que descartás el evento, decís, esto no es un evento, no es un incidente, digo, ya está, listo, lo descartaste. Pero la alerta la tenés que analizar. Sí, ese trabajo te lo tenés que tomar. Entonces, cada evento que llega, cada alerta que llega, se analiza y se destipula, no sé, un acceso, como hablábamos antes, de una persona que está ejecutando una tarea de administración de base de datos y no está dentro de los usuarios permitidos. Eso es un evento, ¿sí? Y yo voy a empezar a investigar cómo esa persona obtuvo sus accesos. quién es ese usuario, si es un usuario que está en nómina o no está en nómina, se empieza a hacer todo ese traqueo de información, ver qué otros eventos se ejecutaron con ese usuario. Voy a empezar a hacer todo ese análisis y va a depender un poco, tenemos ciertos runbooks por cada tipo de evento, qué acciones ir tomando.

    \item[\textbf{MUNTAABSKI FEDERICO:}] Pero un poco es eso, perdón, pero me fui y no me acuerdo exactamente qué era que me habías consultado. No, igual está bien. O sea, está bien que te splashes, pues no, nos sirve también, pues no, nos da información. La pregunta original era, cuando se detectan los eventos, como dijiste vos, en una sesión de un usuario, ¿cómo es el proceso desde la detección hasta

    \item[\textbf{JACUBOVICH LUCIANA:}] por así decirlo, el cierre del evento o la finalización que se dice, esto no es, o sea, no es nada peligroso o esto puede ser un problema de seguridad para la organización. Bueno, un poco es esto que les venía diciendo, ¿no? O sea, es detectar ese evento, poder hacer el análisis de acuerdo al runbook del tipo de evento que sea, o si es algo que está por fuera de los runbooks, poder... o sea, el equipo ya está reunido analizando el evento, y una vez que se determina si hay algún impacto o no, si es un incidente o no, todo eso queda registrado en un informe, que después está abierto a que todo el equipo lo comparta, se sube toda la evidencia registrada durante el análisis, todo esto se... Por más que no sea un incidente, nosotros vamos tomando evidencia para poder tenerla también a futuro, por si se repite alguna situación similar, decir, para, esto ya nos pasó, voy y consulto y tengo la evidencia de lo que sucedía anteriormente en un caso similar. Y es eso, se cierra con el informe. Y muchas veces dentro de ese informe lo que se incluye son mejoras, posibles mejoras que podemos tener también dentro, por ejemplo, de la herramienta del CIEM, decir, no, pará, che, me está alertando esto, que en realidad es cualquier cosa, ¿no? Mejoremos. Entonces, ahí haces otra correlación de eventos que sea mejor. O, uy, pará, este usuario ahora sí tenía que tener estos permisos de administrar de base de datos porque... Pasó de rol y no actualizamos en nuestras reglas del CIEM que esta persona ahora tenía que tener estos permisos y esto no es ningún incidente, simplemente se pone como un falso positivo y se avanza, pero se deja registrado como mejora que hay que cambiar las reglas de la consola del CIEM. Perfecto.

    \item[\textbf{MUNTAABSKI FEDERICO:}] Bueno, pasamos a la siguiente pregunta, un poco también relacionada con las mejoras que vos decís, de tipo, che, esto ya no sirve, o quizás el usuario ahora sí necesitaba tener este rol o registro. La siguiente pregunta sería, ¿cuáles son los principales desafíos que hoy en día vos pensás que hay en lo que es? El monitoreo de usuarios. O sea, pueden ser mejoras, por ejemplo, se me ocurre, menores tiempos o algo que se te ocurra por ese estilo. ¿Mejoras en cuanto al monitoreo de los usuarios? Sí. Sí, que ya no son los usuarios los que hacen las acciones.

    \item[\textbf{JACUBOVICH LUCIANA:}] O sea, el mayor desafío hoy en día es esto que hablábamos al principio de la inteligencia artificial, que hoy ya no es el propio usuario el que está haciendo las acciones, sino que capaz es un agente que lo está haciendo por el usuario. Entonces, a veces poder discernir en eso, o que el usuario tenga esa conciencia de esa autonomía que le está dando. Es la parte difícil de poder detectar esa diferencia entre un usuario y un agente. Entonces, ahí yo creo que están las mejoras que uno puede ir haciendo de poder acercarse a diferenciar esto, porque claramente una cosa es que la acción la da un usuario y otra que le está haciendo un agente automatizado. Y por otro lado, de poder detectar también el uso de esos agentes que a veces no tenemos en radar. O sea, directamente los usuarios los generan y nosotros no tenemos conciencia de que esos agentes existen. Entonces, estar detectando esos agentes hoy dentro de la red.

    \item[\textbf{MOCIULSKY SANTIAGO BERNARDO:}] Perfecto, quedó re claro. Y como otra pregunta más orientada al tema del uso de las herramientas, como comentaste antes, usan herramientas para muchos casos, ¿qué es lo que más se valora en una herramienta al momento de ayudar al equipo de seguridad? O sea, ¿qué son las características que más se toman en cuenta para usar una herramienta y no otra? Los costos.

    \item[\textbf{JACUBOVICH LUCIANA:}] No, la realidad es que hoy este tipo de herramientas de seguridad son súper costosas y cuanto más funcionalidades te dé la herramienta por menor costo, la realidad es que es la herramienta que va a ganar dentro del mercado o inclusive la... La interoperabilidad, digamos, el poder adaptarse a trabajar con otras herramientas de otras marcas diferentes. La realidad es que hoy en día es eso. Hoy para las empresas argentinas, por lo menos, es súper costoso cualquier herramienta que uno quiera adquirir, que no sea open source, que sea paga. Es bastante costoso y si no, si vos querés desarrollar algo más open source, utilizar algo más open source, que por lo general requiere mucha más parametrización, más armado, necesitas mucho más tiempo de CTO y necesitas un equipo más grande para eso. Entonces también es costo. Entonces a veces es como, bueno, ¿dónde pongo?

    \item[\textbf{MOCIULSKY SANTIAGO BERNARDO:}] las ganas va a depender de una decisión más corporativa, directivo, que si se quiere del área propiamente dicho. Ok, y con respecto a, ya fuera del tema de costos, que es totalmente entendible en las empresas, el tema de funcionalidades, o sea, ¿qué es lo que se tomaría en cuenta suponiendo el caso de que el costo no es el principal problema?

    \item[\textbf{JACUBOVICH LUCIANA:}] O sea, ¿qué se tomaría en cuenta como ventaja de una herramienta con respecto a otra? Que cumpla la funcionalidad de lo que vos estás buscando en esa herramienta. O sea, básicamente es eso, que te cubra lo que vos estás buscando y capaz hasta que te dé algún valor abrigado de algo que vos no estabas viendo, que podías llegar a, no sé, pensando en una herramienta de monitoreo, como hablábamos antes. Pensar en monitorear algo que ni se te ocurría que podías monitorear. Creo que es eso un poco. Que sea adaptable, que pueda tener interoperabilidad con otras aplicaciones de otras marcas y que te sumen funcionalidades que vos no tenías en cuenta. Ok, genial.

    \item[\textbf{MOCIULSKY SANTIAGO BERNARDO:}] Otra pregunta, aparte de este tema anterior, es ¿qué factores vos considerás que aumentan el riesgo asociado a un usuario? Por ejemplo, si depende de rol o depende de la actividad que realiza. ¿Qué acciones generan más riesgo? O sea, ¿qué considerás que aumenta el riesgo más en la actividad de un usuario? Si se depende por el rol o por la actividad que realiza. No sé si se llega a entender.

    \item[\textbf{MUNTAABSKI FEDERICO:}] Sí, pero no. Un ejemplo, a ver si capaz me clarifica. O sea, la pregunta está enfocada básicamente en vos monitoreas un grupo de usuarios y lo habíamos pensado tipo, vos siendo por ejemplo un usuario con un rol bajo o teniendo otro con un rol alto. ¿Qué cosas aumentan el riesgo en cada uno de esos casos? ¿Es preferible, por ejemplo, que los usuarios con un rol bajo tengan la misma visibilidad que los de un rol alto? Factores así. No sé si me expliqué bien, la verdad.

    \item[\textbf{JACUBOVICH LUCIANA:}] Seguridad, para mí es como un principio, un pilar, si se quiere, es todo lo que es el mínimo privilegio. Siempre a los usuarios basándose en el mismo, inclusive validando siempre identidades, o sea, si bien como el concepto de Zero Trust, si bien no... es muy difícil de implementar al 100\% el cero trust, como tener esa visión de cero trust, y siempre basándose en el mínimo privilegio, o sea, ahí siempre vas a reducir riesgo, o sea, que el usuario tenga, no quiere decir que no pueda haber un acceso basado en roles, por eso me generaba un poco de confusión, porque vos podés tener un... un rol y que ese rol tenga el mínimo privilegio que requiera el usuario. No necesitas salir de lo que son los accesos por roles para llegar a que el usuario tenga el mínimo privilegio y no incrementar un riesgo. Perfecto. No sé si respondí o no respondí la pregunta. Sí, o sea...

    \item[\textbf{MUNTAABSKI FEDERICO:}] La pregunta quedó, o sea, ahora leyéndola no estaba muy bien formulada, así que... No, no, no. O sea, puedo yo no haber interpretado correctamente por eso. No, no, no. No fue un problema tuyo. La verdad que pensamos que íbamos por un lado, pero como no terminó de... No le pudimos sacar jugo a esta pregunta, porque no la pudimos plantear bien. Bueno, si quieren...

    \item[\textbf{MOCIULSKY SANTIAGO BERNARDO:}] Tienen otra idea, tírenla y le respondo. O sea, no hay problema. Pregúntenme sobre esa, capaz. O sobre lo que les dije. Si se les ocurre, no sé. Dame un medio segundo. Les estoy sacando del libreto. A mí lo que me puse a pensar es... O sea, vos dijiste que monitorean constantemente y hacen, por ejemplo, una investigación de los problemas. capaz la pregunta puede ir por ese lado, o sea, ¿qué son las características que definen que primero investigamos este problema antes que otro? Nosotros tenemos estipulado cuáles son las actividades riesgosas, si se quieren, no sé, privilegios altos, yo voy a estar monitoreando cualquier privilegio elevado o usuarios que tienen, en realidad yo no...

    \item[\textbf{JACUBOVICH LUCIANA:}] monitoreo todos los usuarios, pero lo que yo estoy monitoreando es qué usuarios están realizando una actividad con privilegios elevados. Como les decía antes, el ejemplo de los usuarios con base de datos, con permiso de administración de base de datos. Así como base de datos tenés en todos los aspectos, ¿no? Cualquier tipo de administrador. Entonces, esas son las cosas que yo voy a monitorear después. Tengo otros controles también sobre sus usuarios, más allá del monitoreo, ¿sí? Que me van a disminuir mi riesgo. Que otras actividades riesgosas monitoreo, actividades sospechosas que no son habituales, ¿sí? Entonces, todas esas van a ser alertadas, no sé. la ejecución de un programa que normalmente no se ejecuta. Hay un nuevo programa que se está ejecutando, eso va a ser una alerta. No quiere decir que después termine siendo un incidente, pero sí va a ser alertado. Entonces actividades sospechosas en ese sentido también. Entonces son como diferentes tipos de alertas, de eventos que uno va a monitorear. ¿Sí? Va a dar monitoreos. Hay muchos. Alertas son una porción de eso. ¿Sí? Yo monitoreo, monitoreo el login, el logout, el login fallido. O sea, hay millones de eventos que se monitorean y va a depender también de qué tecnología o qué herramienta, qué aplicación estás monitoreando. Como les dije antes, también en ciertas aplicaciones, ¿cómo se llama? Acciones críticas que se realicen dentro de una aplicación, también los vas a estar monitoreando. Una actividad crítica de cualquier usuario dentro de ese aplicativo, vas a estar monitoreándolo. Entonces, va a depender de qué es lo que estás monitoreando, qué es lo que vas a monitorear. Ahora, de todos esos eventos, que es una bolsa de gatos enorme, es decir, bueno, de todo esto que es lo que estoy alertando, bueno, voy a alertar, sí, lo que me genere más riesgo, o cuando se me dé cierta correlación de esos eventos. Yo siempre pongo el ejemplo más pavo, el que ya no aplica, pero es que cuando antes uno venía todos los días a la oficina y decías, bueno, la máquina la tenías en la oficina, Y no estaba el evento de que la persona había pasado por el molinete y después la persona se conectaba a la máquina y eso tenía que ser un rojo en todos lados, saltando por todos lados, porque cómo puede ser que esa persona se estuviera conectando si no habían ingresado a la oficina. Algún control pasó, algo está fallando. Esos son los tipos de correlaciones de eventos que tenemos que ver. No, esa obviamente hoy en día ya no aplica. Se traslada de otra forma, ¿sí? Por decirlo, si hay una aplicación que es interna, que no está expuesta a internet, si el usuario no se conectó a la VPN, no se puede estar conectando a esa aplicación. Si tenemos todo dentro de esa red securizada, tiene que primero conectarse a la VPN. el usuario entonces tiene que estar ese evento de login del usuario en la vpn para después poder estar ese evento de login en el aplicativo si eso no sucede no sea de esa forma tiene que haber un alerta es como el paralelismo a lo que les decía antes del molinete entonces todo ese tipo de acciones sospechosas y eventos críticos son los que se alertan no es que es Si sucede este evento, se alerta. A veces es, sucede este evento, más este evento, más este evento, o no sucede un evento previo a ese. Si tengo que ir a buscar si previamente sucedió el otro evento, y si no, te lo alerto. Todo eso es, nada, horas de configuración del CIEM, y bueno, habrá mucho de inteligencia artificial también meterle ahí.

    \item[\textbf{MUNTAABSKI FEDERICO:}] cosas de la gente o cosas como las que están pensando en desarrollar ustedes también es parte puede ser parte puede ser parte bueno perfecto ya vamos a ir dándole un cierre a esto ya hicimos creo que la mayoría de las preguntas que tengan muchísimo valor Así que nada, vamos, la última es, básicamente, si hay alguna pregunta que vos pensás que deberíamos haber hecho, que ahora se te ocurra, o que durante la entrevista habrás dicho, che, me tendrían que preguntar esto que no te hicimos. Ok. No, ahí lo que creo que es importante para que tengan, no sé si como preguntas, sino como darles como...

    \item[\textbf{JACUBOVICH LUCIANA:}] información, digamos, que no agrupar a los usuarios bajo una misma categoría, bajo un mismo grupo, por experiencia lo digo. Es decir, bueno, todos los usuarios de finanzas, ustedes hablaban de niveles de riesgo y, o sea, que no, por ser de finanzas tienen un nivel de riesgo y por ser, no, cada usuario es un mundo. O sea, porque por ser de un grupo determinado o tener un rol, que hablábamos antes de los roles, o sea, no. Cada usuario tiene que tener su propio índice de riesgo y que se tiene que ir autogenerando. Y poder cambiar solo. Esa es la parte difícil, si se quiere. Como poder, generando su nuevo rating, ¿sí? Dentro de ese ranking de riesgo, automáticamente en base a su comportamiento, ¿no? Es decir, bueno, está accediendo a esta aplicación, a esta página que es riesgosa o no es riesgosa, como poder hacer ese análisis individualmente por usuario y no como ya presetearle algo por ser de este área o de este otro área. como que tiene que tener ese aprendizaje, la herramienta del usuario. Obviamente que ahí está lo rico que ustedes le tienen que dar de qué es lo riesgoso y qué es lo que no. Pero bueno, ahí entra el desafío. Sí, sí, total. Justo lo que decís es exactamente el desafío que hay es crear ese...

    \item[\textbf{MUNTAABSKI FEDERICO:}] esa línea base de comportamiento, o sea, algo general, y después decir, che, para los usuarios se comportan de forma diferente. Dos gerentes pueden comportarse de forma diferente, pueden hacer actividades totalmente diferentes y no necesariamente es posible ponerlos a todos en una misma bolsa, por así decirlo. Pero bueno, eso serían todas las preguntas.

    \item[\textbf{JACUBOVICH LUCIANA:}] ahí podés parar de grabarse antes si querés ok bueno, buenísimo chicos, no sé lo que necesiten si después se les queda alguna duda o algo me avisan y espero que sea con suerte que salga adelante no se dejen intimidar por Hugo no, no, no Hugo no quiere Bueno, muchísimas gracias por el tiempo y por todo. Esto es invaluable. Bueno, espero que haya servido y que... Nada, si se me ocurre a alguien más que puedan llegar a entrevistar, les aviso. En este momento no... Porque las personas que están en la facu no tienen, o sea, como esta orientación. que yo conozca no quiere decir que no haya otros pero les aviso muchas gracias por todo muchas gracias a ustedes un beso suerte
\end{description}
